\section{Introdução e Motivação}

\begin{frame}{Roteiro}
  \begin{columns}
    \begin{column}{0.5\textwidth}
      \begin{itemize}
        \item \textbf{Parte 1: Contextualização}
        \begin{itemize}
          \item Motivação e Problema
          \item Trabalhos Relacionados
          \item Objetivos
        \end{itemize}
        \vspace{0.3cm}
        \item \textbf{Parte 2: Metodologia}
        \begin{itemize}
          \item Formulação do Problema
          \item Abordagem Proposta
          \item Detalhes Técnicos
        \end{itemize}
      \end{itemize}
    \end{column}
    \begin{column}{0.5\textwidth}
      \begin{itemize}
        \item \textbf{Parte 3: Resultados}
        \begin{itemize}
          \item Experimentos Realizados
          \item Análise dos Resultados
          \item Discussão
        \end{itemize}
        \vspace{0.3cm}
        \item \textbf{Parte 4: Considerações Finais}
        \begin{itemize}
          \item Limitações
          \item Trabalhos Futuros
          \item Conclusão
        \end{itemize}
      \end{itemize}
    \end{column}
  \end{columns}
\end{frame}

\subsection{Motivação e Contexto}
\begin{frame}{Motivação e Contexto}
  \begin{block}{Contexto do Problema}
    Descreva aqui o contexto geral do problema que você está abordando. 
    Qual é o cenário real ou teórico que motiva sua pesquisa?
  \end{block}
  
  \begin{itemize}
    \item \textbf{Desafio principal}: Descreva o desafio central enfrentado.
    \item \textbf{Relevância}: Explique por que esse problema é relevante para a área.
    \item \textbf{Pergunta de pesquisa}:
  \end{itemize}

  \begin{exampleblock}{Pergunta Central}
    Qual é a pergunta de pesquisa que guia este trabalho?
  \end{exampleblock}
\end{frame}

\subsection{Trabalhos Relacionados}
\begin{frame}{Trabalhos Relacionados}
  A abordagem proposta é fundamentada nos seguintes estudos:
  
  \begin{itemize}
    \setlength{\itemsep}{15pt}
    \item \textbf{Autor et al. (Ano)}:
      Breve descrição da contribuição do primeiro trabalho de referência.
    \item \textbf{Autor et al. (Ano)}:
      Breve descrição da contribuição do segundo trabalho de referência.
  \end{itemize}
\end{frame}

\subsection{Objetivos}
\begin{frame}{Objetivos do Trabalho}
  \begin{block}{Objetivo Geral}
    Descreva o objetivo geral do seu trabalho em uma ou duas frases.
  \end{block}
  
  \begin{itemize}
    \setlength{\itemsep}{8pt}
    \item \textbf{Objetivo Específico 1}: Descrição do primeiro objetivo específico.
    \item \textbf{Objetivo Específico 2}: Descrição do segundo objetivo específico.
    \item \textbf{Objetivo Específico 3}: Descrição do terceiro objetivo específico.
  \end{itemize}
\end{frame}
